\documentclass[../../../uem_mat_mei_q.tex]{subfiles}

\begin{document}
    次の空欄中\textsf{ア}，\textsf{ウ}～\textsf{ク}に当てはまるものを選択肢の中から選び%
    その記号をマークせよ．空欄\textsf{イ}には最も適切なものを選択肢の中から選び%
    その記号をマークせよ．それ以外の空欄に当てはまる0から9までの数字を解答用紙の所定の欄に%
    マークせよ．

    \vspace{.5\baselineskip}
    $n$を2以上の自然数とする．原点がOである座標平面を考え，$C$を原点中心で半径が1の円とする．%
    \vspace{1pt}%
    この平面上の点AとPの座標をそれぞれ$(0,\;1)$と$\left(\myfracb{1}{n},\;0\right)$とおく．%
    点Qは次の条件をすべて満たすような点であるとする．
    
    \begin{myenua}
        \item 点Qは円$C$上にある．
            \label{uem_mat_mei_2025_uni1_03_q:1}
        \item $\overrightarrow{\mathrm{AQ}}\cdot\overrightarrow{\mathrm{PQ}}=0$
            \label{uem_mat_mei_2025_uni1_03_q:2}
        \item 点QとAの座標は異なる．
            \label{uem_mat_mei_2025_uni1_03_q:3}
    \end{myenua}
    以下の問いに答えよ．
    \\    
    \begin{minipage}{\linewidth}
        \vspace{.5\baselineskip}
        \centering
            \includegraphics%
            {\subfix{fig_uem_mat_mei_2025_uni1_03/fig_uem_mat_mei_2025_uni1_03_01_q.pdf}}
    \end{minipage}
    \vspace{.1\baselineskip}

    \begin{enumerate}
        \item 線分APの長さは $\myfracd{\sqrt{\,\myboxd{ア}}\,}{n}$ である．
        \item 点Qの座標を$(x,\;y)$とおいて，条件\ref{uem_mat_mei_2025_uni1_03_q:1}，%
            \ref{uem_mat_mei_2025_uni1_03_q:2}を$n$，$x$，$y$の式で表すことにより，%
            $\text{\myboxd{イ}}=0$を得る．さらに条件\ref{uem_mat_mei_2025_uni1_03_q:3}%
            を考慮することにより，点Qの座標は%
            $\left(\myfraca{\myboxd{ウ}}{\myboxa{ア}},%
            \;\myfraca{\myboxd{エ}}{\myboxa{ア}}\right)$%
            であることがわかる．
        \item 3つの点O，P，Aを通る円を$C_1$とすると，$C_1$は点Qも通る．%
            このことに注意すると，$\sin\angle\mathrm{GAP}=%
            \myfraca{\myboxd{オ}}{\myboxc{ア}}$ である．
        \newpage
        \item 三角形PAQについて\\
            $\mathrm{AP}:\mathrm{AQ}:\mathrm{PQ}=%
            \,\myboxc{ア}\,:\,\myboxd{カ}\,:\,\myboxd{キ}$ \:%
            である．
        \item 三角形PAQの内接円の半径を$r$とすると，\\
            $\mathrm{AP}:r=\,\myboxc{ア}\,:\,\myboxd{ク}$ \:である．
        \item $n=2$ とする．点Iを三角形PAQの内接円の中心であるとしたとき

            \myspaceb%
            $\displaystyle%
                    \overrightarrow{\mathrm{QI}}=%
                    \myfraca{1}{\myboxa{ケ}}\overrightarrow{\mathrm{QA}}+%
                    \myfraca{1}{\myboxa{コ}}\overrightarrow{\mathrm{QP}}%
            $\\
            であり，$\overrightarrow{\mathrm{OI}}=%
            \overrightarrow{\mathrm{OQ}}+\overrightarrow{\mathrm{QI}}$ %
            であることから，点Iの座標は%
            $\raisebox{-3pt}{$\Biggl($}\myfraca{1}{\myboxa{サ}},%
            \;\myfraca{1}{\myboxa{シ}}\raisebox{-3pt}{$\Biggr)$}$%
            である．
    
        \vspace{.7\baselineskip}
        \textsf{ア，ウ～クの解答群}
        
        {\renewcommand{\arraystretch}{1.2}%
            \vspace{-.2\baselineskip}%
            \hspace{\zw}%
            \begin{tabularx}{\dimexpr\linewidth-2\zw}{@{}LLLL@{}}
                ⓪\; $0$ & ①\; $1$ & ②\; $n$ & ③\; $2n$ \\
                ④\; $n^2$ & ⑤\; $2n^2$ & ⑥\; $n+1$ & ⑦\; $n-1$ \\
                ⑧\; $n^2+1$ & ⑨\; $n^2-1$
            \end{tabularx}%
        }

        \vspace{.7\baselineskip}
        \textsf{イの解答群}
        
        {\renewcommand{\arraystretch}{1.6}%
            \vspace{-.2\baselineskip}%
            \hspace{\zw}%
            \begin{tabularx}{\dimexpr\linewidth-2\zw}{@{}LLLL@{}}
                ⓪\; $x+\myfracd{y}{n}+1$ & ①\; $x-\myfracd{y}{n}+1$ & ②\; $x+\myfracd{y}{n}-1$ \\
                ③\; $x-\myfracd{y}{n}-1$ & ④\; $\myfracb{x}{n}+y+1$ & ⑤\; $\myfracb{x}{n}-y+1$ \\
                ⑥\; $\myfracb{x}{n}+y-1$ & ⑦\; $\myfracb{x}{n}-y-1$ & ⑧\; $x-ny+1$ \\
                ⑨\; $x+ny+1$
            \end{tabularx}%
        }
    \end{enumerate}
\end{document}
